\chapter{Présentation des entreprises}
\label{chap:entreprises}

% =====================================================================
%  RAPPEL INSTRUCTION : cette partie ne doit EN AUCUN CAS dépasser
%  deux pages au total (les deux entreprises comprises). Rester factuel,
%  éviter le copié-collé de plaquette commerciale.
% =====================================================================

\section{Free Mobile}
Free Mobile est un opérateur de télécommunications du groupe Iliad, fondé
par Xavier Niel. L'entreprise est notamment connue pour être entrée sur le
marché de la téléphonie mobile de manière agressive, en particulier avec son
forfait à deux euros. Au-delà de ses offres commerciales, Free Mobile fait
partie des rares opérateurs à maîtriser l'ensemble de leurs logiciels en
interne : tout leur cœur de réseau y est développé, de même que la
valorisation, et l'entreprise dispose en outre d'une section R\&D dédiée à
ses propres équipements matériels.

J'ai intégré l'équipe Valo, rattachée au Core SI. Cette équipe a la charge
du logiciel de valorisation, qui calcule en temps réel la valeur de l'ensemble des appels
passés en téléphonie mobile sur les différents réseaux.

À mon arrivée, ce logiciel se trouvait dans un contexte de dette technique
élevée : l'ingénieur qui en assurait la maintenance avait quitté
l'entreprise, et plusieurs règles de génie logiciel n'avaient pas été
respectées : une partie des sources, par exemple, n'était pas versionnée. L'équipe, montée et dirigée par David Desmarest, s'est vu confier une
mission de reprise de contrôle de ce logiciel.

Mon rôle y était transverse au travail de l'équipe de développement :
supporter celle-ci en produisant des outils internes, au premier rang
desquels le framework de test détaillé à la section~\ref{chap:t3}.

\section{Beemo Technologie}
Beemo Technologie est une entreprise d'une vingtaine d'années spécialisée
dans la sauvegarde de données. Sa particularité est de ne pas se limiter à
un logiciel installé sur une machine du client, mais de proposer une
solution clé en main : des boîtiers (\emph{box}) déployés au sein de
l'entreprise, voire chez son dirigeant, qui assurent la sauvegarde et la
duplication des données, et donc leur redondance en cas de perte.
Une partie de la solution est également déployée dans le cloud, afin de renforcer cette redondance.

L'entreprise compte une trentaine d'employés et dispose de clients et de
partenaires principalement en France, et s'étend à d'autres pays européens et
africains ; son développement actuel porte sur le continent africain. Une partie des effectifs est dédiée à une \emph{hotline},
reconnue pour sa réactivité et son efficacité dans la résolution des
problèmes. Beemo se distingue enfin par le caractère souverain de sa solution
de sauvegarde, un atout particulièrement recherché dans le contexte actuel.

Mon arrivée chez Beemo a fait suite à une proposition de David Desmarest, qui,
après avoir quitté Free Mobile, y avait pris le \emph{lead} de l'équipe de
développement. Ma décision de le
suivre a répondu à plusieurs facteurs. L'équipe que je quittais traversait une
restructuration qui faisait peser une incertitude sur son devenir, sur laquelle
le management communiquait peu ; je m'y retrouvais paradoxalement parmi les
développeurs les plus anciens, la plupart d'entre eux ayant rejoint l'équipe
après moi. L'orientation méthodologique impulsée par la direction privilégiait
en outre des correctifs rapides, une manière d'aborder la dette technique qui
correspondait moins à la mienne. La proposition de David offrait à l'inverse un
projet dont je connaissais la direction et les perspectives, lui aussi centré
sur la réduction de la dette technique, et la possibilité de revenir à terme au
développement C++ : une perspective qui avait compté dans mon choix initial et
que la situation ne garantissait plus.

À la différence de Free Mobile, où mon rôle était transverse et orienté support,
j'occupais désormais un poste de développeur travaillant directement sur le
fonctionnel. Le contexte restait toutefois marqué par la dette technique :
l'enjeu était de demeurer productif face au marché, c'est-à-dire de
continuer à livrer de nouvelles fonctionnalités, tout en résorbant
progressivement cette dette.

\section{Positionnement des missions}
Sept missions jalonnent ces deux années : cinq chez Free Mobile, deux chez
Beemo Technologie, ces dernières étant encore en cours au moment de la
rédaction. Plutôt que de les présenter dans l'ordre chronologique, restitué
par le diagramme de Gantt en annexe~\ref{chap:annexes}, ce mémoire les
regroupe par nature, en trois chapitres thématiques : \textbf{outiller la productivité et
standardiser l'environnement de développement}, \textbf{rendre testable un
logiciel legacy} pour résorber les régressions et s'en prémunir, et
\textbf{développer du code à visée de production dans un existant endetté}.
Le tableau~\ref{tab:missions} en donne la vue d'ensemble.

\begin{table}[H]
\centering
\begin{tabularx}{\textwidth}{@{}l l X@{}}
\toprule
\textbf{Thème / Entreprise} & \textbf{Mission} & \textbf{Objectif au regard de la dette technique} \\
\midrule
\multicolumn{3}{@{}l}{\emph{Outils : productivité et standardisation de l'environnement de développement}} \\
Free Mobile & Duplication Grafana & Fiabiliser la réplication de la supervision \\
Free Mobile & Compilation C++ & Maîtriser le coût de compilation (CPU/RAM) \\
Beemo      & Rétro-ing. SDK    & Débloquer et documenter une compilation existante \\
\midrule
\multicolumn{3}{@{}l}{\emph{Testabilité d'un logiciel legacy}} \\
Free Mobile & Framework de test & Industrialiser la non-régression \\
Free Mobile & Extension VS Code & Ouvrir l'usage du framework à toute l'équipe \\
Free Mobile & Portage multi-forks & Mutualiser la qualité sur les variantes \\
\midrule
\multicolumn{3}{@{}l}{\emph{Code à visée de production dans un existant endetté}} \\
Beemo      & Immuabilité       & Livrer une fonctionnalité sans accroître la dette \\
\bottomrule
\end{tabularx}
\caption{Vue d'ensemble des missions, regroupées par nature.}
\label{tab:missions}
\end{table}
